\theory{Homology theory}{How can a space's holes be turned into algebraic objects that can be calculated?}{Homology replaces geometric pieces by a chain complex. Cycles are closed pieces, boundaries are pieces that bound something, and homology records cycles that are not boundaries.}{Topology, geometry, data analysis, robotics, mesh computation, algebraic geometry, and differential equations.}{Cycles modulo boundaries measure holes in a way that is invariant under continuous deformation and computable from finite complexes.}

\paragraph{History and the central idea}
Homology grew from Euler's polyhedron formula, Riemann's genus, and Betti's numerical invariants. It became a systematic part of algebraic topology when mathematicians learned to associate algebraic groups with spaces. The word is also used more generally for the homology of any chain complex and for the collection of procedures called homology theories \cite{homology_theory_ref1}.

Think of a triangulated surface. A vertex is a zero-dimensional piece, an edge is a one-dimensional piece, and a triangle is a two-dimensional piece. A collection of edges can form a closed loop. If the loop is the boundary of a collection of triangles, it is not a genuine hole. If it cannot be filled, it represents a one-dimensional hole.

\end{historylist}
\section*{3. Theory}
\subsection*{Core theory}
\paragraph{Chains, cycles, and boundaries}
A chain complex is a sequence of abelian groups or modules
\[
\cdots\longrightarrow C_{n+1}\xrightarrow{\partial_{n+1}}C_n
\xrightarrow{\partial_n}C_{n-1}\longrightarrow\cdots
\]
with $\partial_n\partial_{n+1}=0$. The boundary of a boundary is empty. The $n$-cycles are $Z_n=\ker\partial_n$, and the $n$-boundaries are $B_n=\operatorname{im}\partial_{n+1}$. Since every boundary is a cycle, the $n$th homology group is
\[
H_n=\frac{Z_n}{B_n}.
\]
It identifies cycles that differ by a boundary.

For a connected space, $H_0$ is usually $\mathbb Z$, one copy for each connected component. $H_1$ records loops, and higher groups record higher-dimensional voids. The rank of a homology group is a Betti number; finite pieces such as $\mathbb Z_2$ are torsion.

\paragraph{Examples that build intuition}
The sphere has no one-dimensional hole, so $H_1(S^2)=0$, but it has a two-dimensional surface class, so $H_2(S^2)=\mathbb Z$. A circle has one loop and no filled two-dimensional region: $H_0(S^1)=\mathbb Z$, $H_1(S^1)=\mathbb Z$, and higher homology is zero.

The torus has two independent one-dimensional loops, one around the tube and one through the hole, and one two-dimensional surface class. Its Betti numbers are $1,2,1$. The projective plane has $H_1=\mathbb Z_2$: its essential loop becomes a boundary when traversed twice. This is torsion, a phenomenon not visible from ordinary real-valued measurements \cite{homology_theory_ref2}.

\paragraph{Simplicial and singular homology}
In simplicial homology, a space is decomposed into vertices, edges, triangles, and higher simplices. Boundary matrices make the calculation finite and algorithmic. Singular homology instead allows every continuous map from a standard simplex into the space. It works for arbitrary topological spaces and is invariant under homotopy.

Cellular homology uses cells in a CW complex and is often smaller than singular homology while producing the same answer. Cubical homology uses squares and cubes, which is useful for images and voxel data. Relative homology $H_n(X,A)$ studies chains in $X$ modulo chains already contained in a subspace $A$.

\paragraph{Maps, exact sequences, and products}
A continuous map induces homomorphisms on homology. Homotopic maps induce the same homomorphism, so homology ignores continuous deformation. A pair $A\subseteq X$ gives a long exact sequence linking $H_n(A)$, $H_n(X)$, and relative homology. Mayer--Vietoris calculates a space by splitting it into overlapping pieces.

The cross product combines classes in two spaces, and the Kunneth theorem describes the homology of a product. Poincare duality relates homology and cohomology of an oriented closed manifold. The Euler characteristic is the alternating sum $\chi=\sum(-1)^n\operatorname{rank}H_n$ and is independent of the chosen triangulation.

\paragraph{Homology theories}
Different constructions share the same formal behavior: homotopy invariance, exact sequences, excision, and a dimension rule. Singular, simplicial, cellular, de Rham, and Borel--Moore homology are examples. Generalized theories such as K-theory and cobordism homology relax the ordinary dimension rule and reveal finer structure.

In algebraic geometry, sheaf and coherent homology measure the failure of local data to fit globally. In computational topology, persistent homology studies how homology changes as a scale parameter grows, producing a barcode that identifies robust features in noisy data \cite{homology_theory_ref3}.

\paragraph{Computation and applications}
For a finite mesh, boundary operators are matrices. Row reduction computes ranks and null spaces, while specialized algorithms preserve sparsity. Finite-element software uses homology and cohomology bases to represent fields on meshes. In data analysis, point clouds sampled from a shape are connected at a chosen distance; homology detects components, loops, and voids as the distance changes.

Homology does not reconstruct every detail of a space. Two spaces can have identical homology groups but different fundamental groups or geometry. It is a robust summary, not a complete photograph.

\paragraph{What the theory depends on and where it is useful}
Homology depends on linear algebra, abelian groups, chain complexes, and topology. It helps homotopy theory by providing computable invariants, helps geometry through Euler characteristics and duality, and helps data science by converting shape into persistence diagrams.

\section*{4. Important theorems and results}
\begin{enumerate}[leftmargin=*,itemsep=0.35em]
\item \textbf{Homotopy invariance.} Homotopic maps induce the same map on homology.

\item \textbf{Mayer--Vietoris sequence.} The homology of a union can be computed from the homology of two pieces and their intersection.

\item \textbf{Excision.} Under suitable inclusion conditions, removing a region from a pair does not change relative homology.

\item \textbf{Kunneth theorem.} It describes the homology of a product from the homology of its factors, including possible tensor and torsion terms.

\item \textbf{Poincare duality.} Homology and cohomology groups of an oriented closed manifold are paired in complementary dimensions \cite{homology_theory_ref4}.
\end{enumerate}

\paragraph{Applications}
\begin{enumerate}[leftmargin=*,itemsep=0.35em]
\item \textbf{Topological data analysis and persistent homology.} Persistent homology applies homology to finite point clouds sampled from unknown shapes. The data are embedded in $\mathbb R^d$, and a filtration of simplicial complexes is built by growing balls of radius $\epsilon$ around each point. As $\epsilon$ increases, new connected components appear and merge, new loops form and get filled in, and voids open and close. A persistence diagram or barcode records the birth and death parameter values for each topological feature. Long-lived bars correspond to robust structural features of the underlying space, while short bars are attributed to sampling noise. This pipeline has been applied to reconstruct the shape of proteins from molecular dynamics simulations, to detect loop structures in neuronal connectivity networks, and to identify clustering structure in high-dimensional gene-expression data. The stability theorem for persistence diagrams guarantees that small perturbations of the input point cloud produce only small displacements of the diagram in the bottleneck metric, giving the method robustness that purely algebraic invariants lack \cite{homology_theory_ref3}.
\item \textbf{Sensor networks and coverage analysis.} In a wireless sensor network, each sensor monitors a region modeled as a disk. The union of covered disks forms a topological space, and its homology groups determine whether the network has coverage gaps (one-dimensional holes) or redundant overlaps (higher-dimensional voids). When sensors have unknown locations, pairwise distance measurements construct a \v{C}ech or Rips complex whose homology is computed from boundary matrices. A one-dimensional homology class that persists across a range of distance scales signals a genuine hole in coverage that must be plugged, while the appearance and disappearance of classes at a single scale indicates a transient measurement error. This homology-based approach has been implemented in marine sensor arrays for ocean temperature monitoring, where persistent $H_1$ classes identify regions with insufficient acoustic coverage, and in agricultural IoT networks, where $H_2$ classes detect three-dimensional voids in soil-moisture sensing grids \cite{homology_theory_ref3,homology_theory_ref4}.
\end{enumerate}
\theoryconnections{homology theory}{%
\item Topology supplies spaces and continuous maps, while linear algebra turns chains into boundary matrices and homology groups into kernels modulo images.
\item Category theory explains functoriality, and algebraic geometry and differential geometry provide rich examples.
}{%
\item Homology helps distinguish spaces, compute holes, and prove fixed-point and duality results.
\item It also supplies computable invariants for geometry and topology, making qualitative features accessible through finite calculations in many examples.
}
\section*{Glossary}
\begin{description}[leftmargin=*,style=nextline,itemsep=0.3em]
\item[\textbf{Chain complex.}] A sequence of abelian groups connected by boundary maps whose composition is zero, encoding the combinatorial skeleton of a space.
\item[\textbf{Cycle.}] An element of a chain complex that maps to zero under the boundary operator; it represents a closed piece with no boundary.
\item[\textbf{Boundary.}] An element that is the image of some higher-dimensional chain under the boundary operator; every boundary is automatically a cycle.
\item[\textbf{Homology group.}] The quotient group of cycles modulo boundaries at a given dimension, measuring the extent to which cycles fail to be boundaries.
\item[\textbf{Betti number.}] The rank (free part) of a homology group; it counts independent holes of a given dimension.
\item[\textbf{Torsion.}] A finite cyclic piece of a homology group, invisible to real-valued measurements, indicating a hole that vanishes after being traversed multiple times.
\item[\textbf{Persistent homology.}] A method that tracks how homology groups change as a scale parameter varies, identifying robust topological features in noisy data.
\item[\textbf{Boundary operator.}] The map $\partial_n\colon C_n\to C_{n-1}$ in a chain complex satisfying $\partial^2=0$.
\item[\textbf{Euler characteristic.}] The alternating sum $\chi=\sum(-1)^n\operatorname{rank}H_n$; a topological invariant independent of triangulation.
\item[\textbf{Künneth theorem.}] Describes the homology of a product space in terms of the homology of its factors.
\item[\textbf{Simplicial complex.}] A combinatorial structure built from vertices, edges, triangles, and higher simplices obeying face-inclusion rules.
\item[\textbf{Simplex.}] A generalization of a triangle to higher dimensions: the convex hull of affinely independent points.
\item[\textbf{CW complex.}] A topological space built by iteratively attaching cells of increasing dimension.
\item[\textbf{Excision.}] The theorem that removing a suitable closed subset does not change relative homology.
\end{description}
\section*{7. Sample Questions}
\emph{Exercise reference:} \cite{homology_theory_ref1}. The cited edition is the source for the exercise set; consult its Exercises section for the official statement and numbering.
\begin{enumerate}[leftmargin=*,itemsep=0.9em]
\item \textbf{Exercise 1.} Compute the simplicial homology groups $H_0$ and $H_1$ of the space obtained by identifying the endpoints of three disjoint line segments to a single point (the wedge $S^1\vee S^1\vee S^1$).

\emph{Solution.} The space is a wedge of three circles. Its simplicial chain complex in low dimensions has $C_1=\mathbb{Z}^3$ (three edges) and $C_0=\mathbb{Z}$ (one vertex). The boundary $\partial_1=0$ since each edge starts and ends at the same vertex. So $H_1=\ker\partial_1/\operatorname{im}\partial_2=\mathbb{Z}^3/0=\mathbb{Z}^3$ and $H_0=C_0/\operatorname{im}\partial_1=\mathbb{Z}/0=\mathbb{Z}$ (since the space is connected). Betti numbers: $b_0=1$, $b_1=3$.

\item \textbf{Exercise 2.} Compute $H_*(S^1\times S^1)$, the homology of the torus $T^2$.

\emph{Solution.} By the K\"{u}nneth theorem, $H_n(T^2)=\bigoplus_{p+q=n}H_p(S^1)\otimes H_q(S^1)$. Since $H_0(S^1)=\mathbb{Z}$, $H_1(S^1)=\mathbb{Z}$, and $H_k(S^1)=0$ for $k\geq 2$: $H_0(T^2)=\mathbb{Z}$, $H_1(T^2)=\mathbb{Z}\oplus\mathbb{Z}=\mathbb{Z}^2$, $H_2(T^2)=\mathbb{Z}$. The Euler characteristic is $1-2+1=0$.

\item \textbf{Exercise 3.} Let $X$ be a graph (1-complex) with $5$ vertices and $7$ edges, connected. Compute $H_0(X)$ and $H_1(X)$.

\emph{Solution.} Since $X$ is connected, $H_0(X)=\mathbb{Z}$. For $H_1$: by Euler's formula for a graph, $\chi(X)=V-E=5-7=-2$. Since $\chi=b_0-b_1$, we get $b_1=b_0-\chi=1-(-2)=3$, so $H_1(X)=\mathbb{Z}^3$.

\item \textbf{Exercise 4.} Show that the real projective plane $\mathbb{RP}^2$ has $H_1(\mathbb{RP}^2;\mathbb{Z})\cong\mathbb{Z}/2\mathbb{Z}$.

\emph{Solution.} Using the CW structure with one 0-cell, one 1-cell, and one 2-cell, the cellular chain complex is $0\to\mathbb{Z}\xrightarrow{2}\mathbb{Z}\xrightarrow{0}\mathbb{Z}\to 0$ (the attaching map of the 2-cell wraps twice around the 1-cell). So $H_1=\ker(0)/\operatorname{im}(2)=\mathbb{Z}/2\mathbb{Z}$, and $H_0=\mathbb{Z}$, $H_2=0$.

\item \textbf{Exercise 5.} Use the Mayer--Vietoris sequence to compute $H_1(S^1)$, where $S^1$ is covered by two open arcs $U$ and $V$ overlapping in two disjoint arcs.

\emph{Solution.} $U$, $V$, and $U\cap V$ are all homotopy equivalent to a point (or disjoint union of two contractible arcs). So $H_0(U)=H_0(V)=\mathbb{Z}$, $H_0(U\cap V)=\mathbb{Z}^2$ (two components), and $H_1(U)=H_1(V)=H_1(U\cap V)=0$. The Mayer--Vietoris sequence gives $\cdots\to H_1(U)\oplus H_1(V)\to H_1(S^1)\to H_0(U\cap V)\xrightarrow{\phi}H_0(U)\oplus H_0(V)\to H_0(S^1)\to 0$. This becomes $0\to H_1(S^1)\to\mathbb{Z}^2\xrightarrow{\phi}\mathbb{Z}^2\to\mathbb{Z}\to 0$. The map $\phi$ sends $(a,b)\mapsto(a+b)-(a+b)=\ldots$ more precisely, the difference map sends a generator of each component of $U\cap V$ to its image in $U$ minus its image in $V$: $\phi(a,b)=(a+b,a+b)$... Actually, $\phi(a,b)=(a+b,-a-b)$, so $\ker\phi\cong\mathbb{Z}$ and $\operatorname{im}\phi\cong\mathbb{Z}$. Therefore $H_1(S^1)\cong\ker\phi\cong\mathbb{Z}$.

\end{enumerate}
\section*{8. References}
\addcontentsline{toc}{section}{References}

\begin{thebibliography}{9}
\bibitem{homology_theory_ref1} Allen Hatcher, \emph{Algebraic Topology}, Cambridge University Press, 2002.
\bibitem{homology_theory_ref2} James R. Munkres, \emph{Elements of Algebraic Topology}, Addison-Wesley, 1984.
\bibitem{homology_theory_ref3} Herbert Edelsbrunner and John Harer, \emph{Computational Topology}, American Mathematical Society, 2010.
\bibitem{homology_theory_ref4} Edwin H. Spanier, \emph{Algebraic Topology}, Springer, 1966.
\end{thebibliography}
